%! Linear Transformations and Matrices — Unit 4, Lesson 4.12 — Guided Notes (Answer Key)
\documentclass[10pt]{article}
\usepackage{precalculus-article}
\usepackage{precalculus-key}   % pulls in -boxes; do NOT also load -boxes
\newcommand{\TallMath}[1]{$\displaystyle #1\rule[-1.4em]{0pt}{3.2em}$}

\begin{document}

\pageheader{Unit 4, Lesson 4.12}{Guided Notes --- Answer Key}
\namedateperiod

\vspace{0.1in}

\begin{objectivebox}
By the end of this lesson, I will be able to:
\begin{itemize}[leftmargin=1.5em, itemsep=3pt]
  \item \ans{Apply a linear transformation matrix to an input vector to find the output vector.}
  \item \ans{Transform multiple vectors simultaneously using a single matrix multiplication.}
\end{itemize}
\end{objectivebox}

\vspace{0.08in}

\begin{vocabbox}
\termblanklong{linear transformation}
\ansline{A function $L:\mathbb{R}^2\to\mathbb{R}^2$ where each output component is a sum of constant multiples of the input components.}

\termblanklong{transformation matrix}
\ansline{The unique $2\times 2$ matrix $A$ such that $L(\vec{v})=A\vec{v}$.}

\termblanklong{input vector}
\ansline{The domain vector $\vec{v}$ passed into the transformation.}

\termblanklong{output vector}
\ansline{The range vector $L(\vec{v})=A\vec{v}$.}
\end{vocabbox}

\vspace{0.08in}

\begin{hookbox}
Computer graphics moves every point on screen by transforming position vectors.
If a triangle has vertices at $(1,0)$, $(0,1)$, and $(-1,0)$, can we scale all three
at once using a single matrix multiplication?

\ans{Yes --- place the three vertices as columns of a $2\times 3$ matrix and multiply by the scaling matrix.}
\end{hookbox}

\vspace{0.08in}

\begin{notesbox}{Section 1 --- What is a Linear Transformation? --- Key}
\small
A \textbf{linear transformation} $L:\mathbb{R}^2\to\mathbb{R}^2$ is a function where each
output component is a \ans{sum of constant multiples} of the input components:
\[
  L\!\left(\begin{bmatrix}x\\y\end{bmatrix}\right)
  = \begin{bmatrix}ax+by\\ cx+dy\end{bmatrix}
\]

\textbf{Key fact:} A linear transformation always maps $\vec{0}$ to
\ans{$\vec{0}$}.
That is, $L(\vec{0}) = $ \ans{$\vec{0}$}.

\vspace{0.08in}
\textbf{Zero-vector test:} If $f(\vec{0}) \ne \vec{0}$, then $f$ \ans{cannot}
be a linear transformation.
\end{notesbox}

\vspace{0.08in}

\begin{notesbox}{Section 2 --- Representing Transformations as Matrices --- Key}
\small
Every linear transformation $L:\mathbb{R}^2\to\mathbb{R}^2$ can be written as
$L(\vec{v}) = A\vec{v}$ for a unique $2\times 2$ matrix $A$:
\[
  A = \begin{bmatrix}a & b\\ c & d\end{bmatrix}
\]

\textbf{Example.}  Let $A = \begin{bmatrix}3 & 1 \\ 0 & 2\end{bmatrix}$
and $\vec{v} = \begin{bmatrix}2\\1\end{bmatrix}$.

\vspace{0.06in}
$L(\vec{v}) = A\vec{v} =
\begin{bmatrix}3 & 1 \\ 0 & 2\end{bmatrix}
\begin{bmatrix}2\\1\end{bmatrix}
= $ \ans{$\begin{bmatrix}7\\2\end{bmatrix}$}
\end{notesbox}

\vspace{0.08in}

\begin{notesbox}{Section 3 --- Transforming Multiple Vectors at Once --- Key}
\small
Arrange $n$ input vectors as columns of a $2\times n$ matrix.
One multiplication gives all $n$ output vectors:
\[
  A\begin{bmatrix}\vec{v}_1 & \vec{v}_2 & \cdots & \vec{v}_n\end{bmatrix}
  = \begin{bmatrix}A\vec{v}_1 & A\vec{v}_2 & \cdots & A\vec{v}_n\end{bmatrix}
\]

\textbf{Example.}  Scale triangle vertices $(1,0)$, $(0,1)$, $(-1,0)$ by
$A = \begin{bmatrix}2&0\\0&2\end{bmatrix}$:

\vspace{0.06in}
$\begin{bmatrix}2&0\\0&2\end{bmatrix}
\begin{bmatrix}1 & 0 & -1\\ 0 & 1 & 0\end{bmatrix}
= $ \ans{$\begin{bmatrix}2&0&-2\\0&2&0\end{bmatrix}$}

\vspace{0.08in}
What does this transformation do geometrically?
\ans{It scales (dilates) every point by a factor of 2 from the origin.}
\end{notesbox}

\vspace{0.08in}

\begin{practicebox}
\small
Let $A = \begin{bmatrix}0 & -1 \\ 1 & 0\end{bmatrix}$.
Apply $L$ to $\vec{v}_1 = \begin{bmatrix}1\\0\end{bmatrix}$ and
$\vec{v}_2 = \begin{bmatrix}0\\1\end{bmatrix}$ using a single matrix product.

\vspace{0.08in}
$\begin{bmatrix}0 & -1 \\ 1 & 0\end{bmatrix}
\begin{bmatrix}1 & 0\\ 0 & 1\end{bmatrix} = $ \ans{$\begin{bmatrix}0&-1\\1&0\end{bmatrix}$}

\vspace{0.08in}
What does this transformation appear to do geometrically?
\ans{It rotates every vector $90°$ counterclockwise around the origin.}
\end{practicebox}

\begin{teachernote}
Section 2 example: $3(2)+1(1)=7$ and $0(2)+2(1)=2$, so the output is $\begin{bmatrix}7\\2\end{bmatrix}$.
Section 3: reinforce that the columns of the output matrix are the transformed vertices.
Practice: after computing, ask students to plot the input and output vectors to see the rotation.
\end{teachernote}

\end{document}
